Mit $v = 2\pi r/T$ ist die Radialbeschleunigung bei gleichem Radius
\begin{align}
 a_\text{r} &= \frac{v^2}{r} = \frac{4\pi^2 r}{T^2} \propto \frac{1}{T^2}
\end{align}
Mit $T' = (1-p)\,T$ und $p = \pkO$ wird
\begin{align}
 \frac{a_\text{r}'}{a_\text{r}} &= \left(\frac{T}{T'}\right)^2 = \faF = \fa
\end{align}
Die Radialbeschleunigung nimmt also zu um
\begin{align}
 \frac{a_\text{r}'-a_\text{r}}{a_\text{r}} &= \daF = \da \approx \result{\daP{0}{2}}
\end{align}
